%! Function Fundamentals — Unit 1, Lesson 1.1 — Slides (Beamer)
\documentclass[aspectratio=169, 11pt]{beamer}
\usepackage{precalculus-beamer}   % provides \plumheader{} and \sectionlabel[color]{}

\begin{document}

% TITLE SLIDE
{
\setbeamercolor{background canvas}{bg=plum}
\begin{frame}[plain]
  \vspace{2.8cm}\hspace{0.55cm}%
  \begin{minipage}{0.9\textwidth}
    {\color{goldacc}\bfseries\small LESSON 1.1}\\[0.5cm]
    {\color{white}\bfseries\LARGE Function Fundamentals}\\[1.2cm]
    {\color{lilacmid}\small Precalculus~$\cdot$~Unit 1: Functions and Change}
  \end{minipage}
\end{frame}
}

% HOOK SLIDE
\begin{frame}[t, plain]
  \plumheader{Hook: The Machine Gets a Name}
  \sectionlabel[goldacc]{Think About It}

  Yesterday's trustworthy vending machine is back --- this time with a label.
  Let \(P(b)\) be the price in dollars of the item at button \(b\).

  \bigskip
  \begin{center}
    {\Huge\bfseries\color{plum} \(P(\text{B4}) = 1.25\)}
  \end{center}

  \bigskip
  \begin{columns}[T]
    \column{0.48\textwidth}
      \begin{block}{Say it out loud}
        What is this equation \textit{saying}? Give a full English sentence, not
        a reading of the symbols.
      \end{block}

    \column{0.48\textwidth}
      \begin{block}{Then ask}
        What is the \textbf{input} here? Does an input have to be a
        \textbf{number}?
      \end{block}
  \end{columns}

  \bigskip
  \textit{One short symbol just packed in the machine, the input, and the
  output. That is the whole point of function notation.}
\end{frame}

% SECTION 1 — FUNCTION NOTATION
\begin{frame}[t, plain]
  \plumheader{1. Function Notation --- Naming the Machine}

  \begin{block}{Three parts, one symbol}
    In \(f(x)\): \quad \(f\) is the \textbf{name} of the function \(\cdot\)
    \(x\) is the \textbf{input} \(\cdot\) \(f(x)\) is the \textbf{output}.
  \end{block}

  \medskip
  Read \(f(x)\) as ``\textbf{\(f\) of \(x\)}.'' It is \textbf{not} \(f\) times
  \(x\) --- the parentheses hold the input, they do not multiply.

  \smallskip
  \begin{columns}[T]
    \column{0.48\textwidth}
      \textbf{In context.} A rideshare fare:
      \[ F(m) = 2.50 + 1.75m \]
      \begin{itemize}\setlength{\itemsep}{5pt}
        \item Name: \(F\)
        \item Input \(m\): miles driven
        \item Output \(F(m)\): the fare, in dollars
        \item So \(F(4)\) asks: what does a 4-mile ride cost?
      \end{itemize}

    \column{0.48\textwidth}
      \textbf{One statement, three facts.} \(f(3) = 7\) says:
      \begin{itemize}\setlength{\itemsep}{5pt}
        \item the input is \textbf{3}
        \item the output is \textbf{7}
        \item the point \(\mathbf{(3, 7)}\) is on the graph of \(f\)
      \end{itemize}
  \end{columns}
\end{frame}

% SECTION 2 — EVALUATING FROM A FORMULA
\begin{frame}[t, plain]
  \plumheader{2. Evaluating from a Formula}

  \begin{block}{The empty-slot habit}
    The slot does not care what you put in it:
    \[ f(x) = 5x - 8 \qquad\longrightarrow\qquad f(\square) = 5(\square) - 8 \]
    Write the parentheses \textbf{first}, empty. Then fill them.
  \end{block}

  \bigskip
  \begin{columns}[T]
    \column{0.31\textwidth}
      \textbf{Find \(f(3)\)}
      \begin{align*}
        f(3) &= 5(3) - 8 \\
             &= 7
      \end{align*}

    \column{0.31\textwidth}
      \textbf{Find \(f(-2)\)}
      \begin{align*}
        f(-2) &= 5(-2) - 8 \\
              &= -18
      \end{align*}

    \column{0.34\textwidth}
      \textbf{\(g(x) = x^2 + 3x\); find \(g(-4)\)}
      \begin{align*}
        g(-4) &= (-4)^2 + 3(-4) \\
              &= 16 - 12 = 4
      \end{align*}
  \end{columns}

  \bigskip
  \textbf{The trap:} \((-4)^2 = 16\), but \(-4^2 = -16\). Evaluating never
  changes the rule --- it just runs the machine once.
\end{frame}

% SECTION 3 — TWO DIRECTIONS
\begin{frame}[t, plain]
  \plumheader{3. Two Directions: Evaluate vs.\ Solve}
  \sectionlabel[redacc]{The Hinge of This Lesson}

  \begin{tabular}{|p{0.22\textwidth}|p{0.32\textwidth}|p{0.34\textwidth}|}
  \hline
   & \textbf{Find \(f(2)\)} & \textbf{Solve \(f(x) = 2\)} \\
  \hline
  You are given & the \textbf{input} & the \textbf{output} \\
  \hline
  You are asked for & the \textbf{output} & the \textbf{input(s)} \\
  \hline
  How many answers & exactly one & zero, one, or many \\
  \hline
  \end{tabular}

  \bigskip
  \textbf{From a table.} Each column is one input--output pair.

  \medskip
  \begin{tabular}{|c|c|c|c|c|c|}
  \hline
  \(x\) & 1 & 2 & 3 & 4 & 5 \\
  \hline
  \(s(x)\) & 12 & 9 & 9 & 6 & 3 \\
  \hline
  \end{tabular}

  \medskip
  \(s(4) = ?\) \quad (read \textit{down}) \qquad Solve \(s(x) = 9\)
  \quad (read \textit{up} --- and it happens twice)
\end{frame}

% SECTION 3, CONT. — THE GRAPH
\begin{frame}[t, plain]
  \plumheader{Reading Both Directions off a Graph}

  \begin{columns}[T]
    \column{0.44\textwidth}
      \begin{center}
      \begin{tikzpicture}[x=0.55cm, y=0.42cm]
        \draw[linegray, very thin, step=1] (0,0) grid (8,7);
        \draw[->, thick] (0,0) -- (8.6,0) node[below right] {\small \(x\)};
        \draw[->, thick] (0,0) -- (0,7.6) node[above] {\small \(f(x)\)};
        \foreach \x in {1,2,3,4,5,6,7,8} \node[below] at (\x,0) {\small \x};
        \foreach \y in {1,3,5,7} \node[left] at (0,\y) {\small \y};
        \draw[very thick, plum] (0,5) -- (2,1);
        \draw[very thick, plum] (2,1) -- (5,7);
        \draw[very thick, plum] (5,7) -- (8,4);
        \fill[plum] (0,5) circle (2.5pt);
        \fill[plum] (8,4) circle (2.5pt);
      \end{tikzpicture}
      \end{center}

    \column{0.52\textwidth}
      \begin{itemize}\setlength{\itemsep}{7pt}
        \item \(f(2) = ?\) --- start at \(x = 2\), walk \textbf{up} to the
              graph, read across
        \item \(f(6) = ?\)
        \item Solve \(f(x) = 3\) --- start at 3 on the \textbf{output} axis and
              walk \textbf{across}: two answers
        \item Solve \(f(x) = 5\) --- \textbf{three} answers
      \end{itemize}
  \end{columns}

  \bigskip
  Three inputs sharing one output is \textbf{many-to-one} --- allowed, and
  \(f\) is still a function. One \textit{input} with two outputs is what breaks
  the rule.
\end{frame}

% SECTION 4 — DOMAIN AND RANGE
\begin{frame}[t, plain]
  \plumheader{4. Domain and Range}

  \begin{block}{Two words for two sets}
    The \textbf{domain} is the set of all \textbf{inputs} the function uses. The
    \textbf{range} is the set of all \textbf{outputs} it produces --- a repeated
    output is listed \textbf{once}.
  \end{block}

  \medskip
  \begin{columns}[T]
    \column{0.48\textwidth}
      \textbf{From a list of points}
      \[ f = \{(-2, 6),\ (0, 4),\ (3, 6),\ (5, 1)\} \]
      Domain: \(\{-2, 0, 3, 5\}\) \\[3pt]
      Range: \(\{1, 4, 6\}\) --- the 6 is written once

      \medskip
      \textbf{From the graph on the last slide}\\[3pt]
      Domain: \(0 \le x \le 8\) \quad Range: \(1 \le y \le 7\)

    \column{0.48\textwidth}
      \begin{block}{Notation upgrade}
        \(0 \le x \le 8\) is also written \([\,0, 8\,]\) --- a square bracket
        includes the endpoint, a round one does not. Both forms count as
        correct.
      \end{block}

      \medskip
      \textbf{From a formula:} \(k(x) = \dfrac{8}{x + 3}\)

      \smallskip
      At \(x = -3\) the denominator is 0 --- \textbf{undefined}. Domain: all
      real numbers \textbf{except} \(-3\).
  \end{columns}
\end{frame}

% CLASS PRACTICE
\begin{frame}[t, plain]
  \plumheader{Class Practice}
  \sectionlabel[redacc]{Try It}

  \begin{enumerate}\setlength{\itemsep}{10pt}
    \item Let \(f(x) = 3x - 7\). Find \(f(5)\) and \(f(-1)\).
    \item The table shows a function \(p\).

          \smallskip
          \begin{tabular}{|c|c|c|c|c|}
          \hline
          \(x\) & 1 & 2 & 3 & 4 \\
          \hline
          \(p(x)\) & 9 & 6 & 6 & 3 \\
          \hline
          \end{tabular}

          \smallskip
          \begin{enumerate}\setlength{\itemsep}{5pt}
            \item Find \(p(2)\). Then solve \(p(x) = 3\). Which one was an
                  \textit{evaluate} question?
            \item State the domain and the range of \(p\).
          \end{enumerate}
  \end{enumerate}
\end{frame}

\end{document}
